\documentclass[12pt,letterpaper]{article}
\usepackage[margin=1in]{geometry}
\usepackage[T1]{fontenc}
\usepackage{amsmath,amssymb}
\usepackage{newtxtext,newtxmath}
\usepackage{helvet}
\usepackage{enumitem}
\usepackage{caption,booktabs}
\usepackage[hidelinks]{hyperref}
\captionsetup{font=small,labelfont=bf,skip=4pt}
\renewcommand{\arraystretch}{1.15}
\setlength{\emergencystretch}{2em}
\pagestyle{empty}
\setlength{\parindent}{0pt}
\setlength{\parskip}{0pt}
\AtBeginDocument{\fontsize{12}{13.8}\selectfont}
\setlist[enumerate,1]{label=\arabic*.,leftmargin=0.25in,
labelsep=0.08in,align=left,topsep=0pt,itemsep=14pt,parsep=6pt}
\newcommand{\um}{\,\mu\mathrm{m}}
\newcommand{\uA}{\,\mu\mathrm{A}}
\newcommand{\mS}{\,\mathrm{mS}}
\newcommand{\kohm}{\,\mathrm{k}\Omega}
\newcommand{\Mohm}{\,\mathrm{M}\Omega}
\newcommand{\V}{\,\mathrm{V}}
\newcommand{\MHz}{\,\mathrm{MHz}}
\newcommand{\Vus}{\,\mathrm{V}/\mu\mathrm{s}}

\hypersetup{pdfauthor={Anonymous contributor},pdftitle={ECEN 474/704 - Lab 6 Prelab},pdfsubject={Anonymized laboratory prelab}}
\begin{document}
{\fontsize{11}{12.65}\selectfont
\hfill October 9, 2026\par}
{\centering
{\fontsize{14}{16.1}\selectfont\sffamily\bfseries ECEN 474/704\par}
{\fontsize{14}{16.1}\selectfont Lab 6: Prelab\par}
{\fontsize{14}{16.1}\selectfont Anonymous contributor\par}
}
\vspace{12pt}
\begin{enumerate}

\item \textbf{Comparison of differential-amplifier topologies.}
Compare the two single-ended differential amplifiers in this lab.
Rate differential gain, common-mode gain, output voltage swing,
CMRR, and common-mode input voltage range. Include expressions
for each specification.

\textbf{Solution.}

The comparison assumes comparable currents and device parameters.
``Good'' denotes larger differential gain, smaller
common-mode gain magnitude, larger voltage range, or higher CMRR.

\begin{center}
\small
\setlength{\tabcolsep}{4pt}
\begin{tabular}{@{}lcc@{}}
\toprule
Specification & Simple (Fig.~6-6) & Cascode (Fig.~6-8)\\
\midrule
Differential gain & Medium & Good\\
Common-mode gain ($|A_{CM}|$) & Medium & Good\\
Output voltage swing & Good & Bad\\
CMRR & Medium & Good\\
Common-mode input voltage range & Good & Bad\\
\bottomrule
\end{tabular}
\end{center}

With $V_{DM}=V_{i+}-V_{i-}$ and
$V_{CM}=(V_{i+}+V_{i-})/2$, zero differential input gives
$I_{D1}=I_{D2}=I_{\mathrm{tail}}/2$.
The overdrive voltages are $V_{OVn}=V_{GS}-V_{TN}$ and
$V_{OVp}=V_{SG}-|V_{TP}|$, both positive.

\textbf{(a) Differential gain.}
The single-ended gain has no factor of $1/2$ because the PMOS
mirror combines the two input-branch signal currents:
\begin{align*}
A_{DM,\mathrm{simple}}&\approx g_{m1}(r_{o2}\parallel r_{o6}),\\
R_{P,\mathrm{cas}}&\approx r_{o8}+r_{o10}+g_{m8}r_{o8}r_{o10},\\
A_{DM,\mathrm{cas}}&\approx
g_{m1}(r_{o2}\parallel R_{P,\mathrm{cas}})\approx g_{m1}r_{o2}.
\end{align*}
The cascode increases the load resistance, but M2 is not cascoded;
its $r_{o2}$ therefore limits the differential gain.

\textbf{(b) Common-mode gain.}
Using the low-frequency approximations in the Lab 6 handout,
\begin{align*}
A_{CM,\mathrm{simple}}&\approx-\frac{1}{2r_{o4}g_{m5}},\\
R_{SS,\mathrm{cas}}&\approx r_{o4}+r_{o6}+g_{m4}r_{o4}r_{o6},\\
A_{CM,\mathrm{cas}}&\approx-\frac{1}{2R_{SS,\mathrm{cas}}g_{m9}}.
\end{align*}
The cascode tail source has a higher resistance and reduces $|A_{CM}|$.

\textbf{(c) Common-mode rejection ratio.}
For either topology,
\[
\boxed{\mathrm{CMRR}=20\log_{10}\!\left|\frac{A_{DM}}{A_{CM}}\right|.}
\]
\newpage
\textbf{(d) Common-mode input voltage range.}
For the simple amplifier, the saturation constraints give
\begin{align*}
V_{CM,\min}&=V_{SS}+V_{OV4}+V_{TN1}+V_{OV1},\\
V_{CM,\max}&=V_{DD}-|V_{TP5}|-V_{OV5}+V_{TN1},\\
\Delta V_{CM}&=V_{CM,\max}-V_{CM,\min}.
\end{align*}
For the cascode amplifier, the voltage reserved across
the lower NMOS mirror device is
$d_n=V_{S4}-V_{SS}$. With a fixed $V_{BN}$,
\begin{align*}
d_n&=V_{BN}-V_{SS}-V_{TN4}-V_{OV4},\\
V_{CM,\min}&=V_{SS}+d_n+V_{OV4}+V_{TN1}+V_{OV1},\\
V_{CM,\max}&=V_{DD}-|V_{TP9}|-V_{OV9}+V_{TN1}.
\end{align*}
The handout's minimum-headroom expression is the special case
$d_n=V_{OV6}$. A practical choice $d_n>V_{OV6}$ provides
saturation margin but raises the lower input limit. These ranges
also require the reference branches and output devices to remain
in saturation; output bias can impose an additional input limit.

\textbf{(e) Output voltage swing.}
Let $V_S$ be the common source voltage of M1 and M2. At a fixed
common-mode input and with matched input devices,
\[
V_{o,\min}=V_S+V_{OV2}=V_{CM}-V_{TN1}.
\]
For the simple circuit, the upper output limit is
$V_{o,\max}=V_{DD}-V_{OV6}$.
For the cascode circuit, $d_p=V_{DD}-V_{S8}$ gives
\[
d_p=V_{DD}-V_{BP}-|V_{TP8}|-V_{OV8},\qquad
V_{o,\max}=V_{DD}-d_p-V_{OV8}.
\]
The ideal choice $d_p=V_{OV10}$ gives
$V_{o,\max}=V_{DD}-V_{OV10}-V_{OV8}$.
In both cases the available interval width is
$\Delta V_o=V_{o,\max}-V_{o,\min}$, while the largest symmetric
signal about $V_{o,Q}$ has peak amplitude
\[
v_{o,\mathrm{pk}}<\min(V_{o,Q}-V_{o,\min},\,V_{o,\max}-V_{o,Q}).
\]
These are saturation-based bounds; a specified distortion limit
would require additional margin inside the interval.

\textbf{Parameters and assumptions for Questions 2 and 3.}
The single-finger $W/L=5/0.27\um$ data in Table I of the
Lab 1 report provide the device parameters:
\begin{center}
\small
\begin{tabular}{@{}lcc@{}}
\toprule
Parameter & NFET & PFET\\
\midrule
$K_P=\mu C_{ox}$ ($\mu$A/V$^2$) & 215.6 & 47.57\\
$|V_{TH}|$ (V) & 0.4330 & 0.4068\\
$\lambda$ (V$^{-1}$) & 0.1624 & 0.1945\\
$\theta$ (V$^{-1}$) & 0.9999 & 0.6558\\
\bottomrule
\end{tabular}
\end{center}
All initial lengths are $0.27\um$. The DC sizing uses the square-law
model and neglects body effect, mobility degradation, and
channel-length-modulation corrections to the drain current.
The parameter $\lambda$ enters only the small-signal estimate
$r_o\approx1/(\lambda I_D)$. All widths are total effective widths.
The resulting dimensions constitute an analytical starting point.

\newpage
\item \textbf{Simple differential-amplifier design (Fig.~6-6).}
Design for $\mathrm{SR}>15\Vus$, $\mathrm{GBW}>8\MHz$, and
$\Delta V_{CM}>0.4\V$, with $V_{DD}=0.9\V$,
$V_{SS}=-0.9\V$, and $C_L=10\,\mathrm{pF}$.

\textbf{Solution.}

\textbf{(a) Tail current and input transconductance.}
The slew-rate constraint requires
\[
I_{\mathrm{tail}}>C_L(15\times10^6)=150\uA.
\]
The handout gives angular bandwidth $\omega_{GBW}=g_{m1}/C_L$.
Converting to hertz gives
\[
\mathrm{GBW}=\frac{g_{m1}}{2\pi C_L},\qquad
g_{m1}>2\pi(8\times10^6)(10\times10^{-12})=0.503\mS.
\]
For $I_{\mathrm{tail}}=I_{\mathrm{bias}}=180\uA$ and a
1:1 M3--M4 mirror, each input and PMOS load device carries
$90\uA$ at the balanced operating point. The selected overdrives are
\[
V_{OV1,2}=0.250\V,\quad V_{OV3,4}=0.150\V,\quad
V_{OV5,6}=0.200\V.
\]
Then $g_{m1}=2(90\uA)/(0.250\V)=0.720\mS$, providing
margin above the required transconductance.

\textbf{(b) Transistor widths and lengths.}
The square-law sizing relations are
\[
I_D\approx\frac12 K_P\frac WL V_{OV}^2,\qquad
\frac WL=\frac{2I_D}{K_PV_{OV}^2},\qquad W=L\frac WL.
\]
The sizing calculations, with currents in $\mu$A, are
\begin{align*}
(W/L)_{1,2}&=\frac{2(90)}{215.6(0.250)^2}=13.358,\\
(W/L)_{3,4}&=\frac{2(180)}{215.6(0.150)^2}=74.212,\\
(W/L)_{5,6}&=\frac{2(90)}{47.57(0.200)^2}=94.597.
\end{align*}
\begin{center}
\small
\begin{tabular}{@{}lcccc@{}}
\toprule
Device & $I_D$ ($\mu$A) & $V_{OV}$ (V) & $W$ ($\mu$m) & $L$ ($\mu$m)\\
\midrule
M1, M2, each & 90 & 0.250 & 3.607 & 0.270\\
M3, M4, each & 180 & 0.150 & 20.037 & 0.270\\
M5, M6, each & 90 & 0.200 & 25.541 & 0.270\\
\bottomrule
\end{tabular}
\end{center}
NMOS bodies connect to $V_{SS}$ and PMOS bodies to $V_{DD}$.
Matched device pairs use equal dimensions and consistent finger
configurations. The diode-connected M3 establishes the common gate
voltage, eliminating the need for a separate tail-bias voltage source.

\textbf{(c) Common-mode range.}
Substituting the selected overdrives gives
\begin{align*}
V_{CM,\min}&=-0.9+0.150+0.433+0.250=-0.067\V,\\
V_{CM,\max}&=0.9-0.4068-0.200+0.433=0.7262\V,\\
\Delta V_{CM}&=\boxed{0.7932\V>0.4\V.}
\end{align*}
The selected bias $V_{CM,Q}=0.300\V$ lies inside this interval.

\newpage
\textbf{(d) Nominal DC operating point and saturation checks.}
For equal inputs, a symmetry-consistent operating point places the
output at the diode-connected PMOS voltage:
\begin{align*}
V_S&=0.300-(0.433+0.250)=-0.3830\V,\\
V_{D1}=V_{o,Q}&=0.900-(0.4068+0.200)=0.2932\V,\\
V_{G3}=V_{G4}&=-0.900+0.433+0.150=-0.3170\V.
\end{align*}
Each device satisfies $V_{DS}>V_{OV}$ for an NMOS or
$V_{SD}>V_{OV}$ for a PMOS under the assumed model:
\begin{center}
\small
\begin{tabular}{@{}lccc@{}}
\toprule
Device & $V_{DS}$ or $V_{SD}$ (V) & $V_{OV}$ (V) & Saturation\\
\midrule
M1, M2 & 0.6762 & 0.250 & Yes\\
M3 & 0.5830 & 0.150 & Yes\\
M4 & 0.5170 & 0.150 & Yes\\
M5, M6 & 0.6068 & 0.200 & Yes\\
\bottomrule
\end{tabular}
\end{center}
At $V_{CM,Q}=0.300\V$, the output limits are
\[
-0.133\V<V_o<0.700\V.
\]
The interval is $0.833\V$ wide. About $V_{o,Q}=0.2932\V$,
the symmetric-swing limit approaches $0.4068\V$ peak. A finite
saturation margin reduces the usable amplitude below this limit.

\textbf{(e) Small-signal performance.}
Using the stated $\lambda$ values,
\begin{align*}
r_{o2}&=\frac{1}{0.1624(90\times10^{-6})}=68.418\kohm,\\
r_{o6}&=\frac{1}{0.1945(90\times10^{-6})}=57.127\kohm,\\
r_{o4}&=\frac{1}{0.1624(180\times10^{-6})}=34.209\kohm,\\
R_{\mathrm{out}}&=r_{o2}\parallel r_{o6}=31.132\kohm.
\end{align*}
Thus $A_{DM}=g_{m1}R_{\mathrm{out}}=22.415$ V/V
($27.01$ dB). With $g_{m5}=2(90\uA)/(0.200\V)=0.900\mS$,
\[
A_{CM}\approx-\frac{1}{2r_{o4}g_{m5}}=-0.01624,\qquad
\mathrm{CMRR}\approx62.80\,\mathrm{dB}.
\]
The dominant-pole approximation gives
\[
f_{p1}=\frac{1}{2\pi R_{\mathrm{out}}C_L}=511.22\,\mathrm{kHz},\quad
\mathrm{GBW}=\frac{g_{m1}}{2\pi C_L}=11.46\MHz.
\]
\begin{center}
\small
\begin{tabular}{@{}lcc@{}}
\toprule
Specification & Requirement & Analytical estimate\\
\midrule
Slew rate, $I_{\mathrm{tail}}/C_L$ & $>15\Vus$ & $18\Vus$\\
GBW & $>8\MHz$ & $11.46\MHz$\\
Common-mode range width & $>0.4\V$ & $0.7932\V$\\
\bottomrule
\end{tabular}
\end{center}
\newpage
\item \textbf{Differential amplifier with cascode mirrors (Fig.~6-8).}
Design for $\mathrm{SR}>7\Vus$, $\mathrm{GBW}>8\MHz$, and
$\mathrm{CMRR}>65$ dB, with $V_{DD}=0.9\V$,
$V_{SS}=-0.9\V$, and $C_L=10\,\mathrm{pF}$.
Determine $V_{BP}$ and $V_{BN}$ and retain a nonzero input
common-mode range.

\textbf{Solution.}

\textbf{(a) Current, overdrives, and transistor dimensions.}
The minimum tail current is $C_L(7\times10^6)=70\uA$;
the input-transconductance minimum remains $0.503\mS$.
A 1:1 tail mirror establishes
\[
I_{\mathrm{bias}}=I_{\mathrm{tail}}=120\uA,\qquad
I_{D1,2,7,8,9,10}=60\uA.
\]
For $V_{OV}=0.150\V$ in all ten devices and $L=0.270\um$,
the sizing equation of Question 2 gives
\begin{align*}
(W/L)_{1,2}&=\frac{2(60)}{215.6(0.150)^2}=24.737,\\
(W/L)_{3,4,5,6}&=\frac{2(120)}{215.6(0.150)^2}=49.474,\\
(W/L)_{7,8,9,10}&=\frac{2(60)}{47.57(0.150)^2}=112.115.
\end{align*}
\begin{center}
\small
\begin{tabular}{@{}lcccc@{}}
\toprule
Device & $I_D$ ($\mu$A) & $V_{OV}$ (V) & $W$ ($\mu$m) & $L$ ($\mu$m)\\
\midrule
M1, M2, each & 60 & 0.150 & 6.679 & 0.270\\
M3--M6, each & 120 & 0.150 & 13.358 & 0.270\\
M7--M10, each & 60 & 0.150 & 30.271 & 0.270\\
\bottomrule
\end{tabular}
\end{center}
The input transconductance is
$g_{m1}=2(60\uA)/(0.150\V)=0.800\mS$.

\textbf{(b) Cascode bias voltages.}
Reserving $d_n=d_p=0.200\V$ across the lower NMOS and upper
PMOS mirror devices gives $50$ mV of nominal saturation
margin above their $0.150$ V overdrives:
\begin{align*}
V_N=V_{S3}=V_{S4}&=V_{SS}+d_n=-0.700\V,\\
V_P=V_{S7}=V_{S8}&=V_{DD}-d_p=0.700\V,\\
V_{BN}&=V_N+V_{TN4}+V_{OV4}=\boxed{-0.117\V,}\\
V_{BP}&=V_P-|V_{TP8}|-V_{OV8}=\boxed{0.1432\V.}
\end{align*}
As permitted by the handout, external sources at the testbench
supply these biases through additional schematic pins. The amplifier
therefore requires no internal bias-generation transistors or
additional bias current for these voltages.

\textbf{(c) Common-mode input range with these biases.}
Including the selected $d_n$, rather than only its ideal minimum,
\begin{align*}
V_{CM,\min}&=-0.900+0.200+0.150+0.433+0.150=0.0330\V,\\
V_{CM,\max}&=0.900-0.4068-0.150+0.433=0.7762\V.
\end{align*}
Hence $\boxed{\Delta V_{CM}=0.7432\V>0}$.
The nominal operating point is $V_{CM,Q}=0.300\V$.

\newpage
\textbf{(d) DC operating point and saturation checks.}
In Fig.~6-8, M3's drain drives the gates of M5 and M6; M7's
drain drives the gates of M9 and M10. Therefore,
\begin{align*}
V_{G5}=V_{G6}=V_{D3}&=-0.900+0.433+0.150=-0.3170\V,\\
V_{D1}=V_{G9}=V_{G10}=V_{o,Q}&=0.900-0.4068-0.150=0.3432\V,\\
V_S&=0.300-0.433-0.150=-0.2830\V.
\end{align*}
Together with $V_N=-0.700\V$ and $V_P=0.700\V$, these
voltages give the following first-order saturation checks:
\begin{center}
\small
\begin{tabular}{@{}lccc@{}}
\toprule
Device & $V_{DS}$ or $V_{SD}$ (V) & $V_{OV}$ (V) & Saturation\\
\midrule
M1, M2 & 0.6262 & 0.150 & Yes\\
M3 & 0.3830 & 0.150 & Yes\\
M4 & 0.4170 & 0.150 & Yes\\
M5, M6 & 0.2000 & 0.150 & Yes\\
M7, M8 & 0.3568 & 0.150 & Yes\\
M9, M10 & 0.2000 & 0.150 & Yes\\
\bottomrule
\end{tabular}
\end{center}
At the chosen common-mode voltage, the output interval is
$-0.133\V<V_o<0.550\V$. Its width is $0.683\V$, and
the symmetric swing about $V_{o,Q}=0.3432\V$ approaches
$0.2068\V$ peak. The smaller upper headroom is the cost of
the cascode PMOS load.

\textbf{(e) Gain, common-mode rejection, and frequency estimates.}
The input and PMOS transconductances are $0.800\mS$;
the NMOS tail-device transconductances are $1.600\mS$.
Using $r_o\approx1/(\lambda I_D)$ gives
\[
r_{o2}=102.627\kohm,\quad
r_{o4}=r_{o6}=51.314\kohm,\quad
r_{o8}=r_{o10}=85.690\kohm.
\]
Consequently,
\begin{align*}
R_{SS}&\approx r_{o4}+r_{o6}+g_{m4}r_{o4}r_{o6}=4.3156\Mohm,\\
R_P&\approx r_{o8}+r_{o10}+g_{m8}r_{o8}r_{o10}=6.0456\Mohm,\\
R_{\mathrm{out}}&\approx r_{o2}\parallel R_P=100.914\kohm,\\
A_{DM}&\approx g_{m1}R_{\mathrm{out}}=80.731\ \mathrm{V/V}
\quad(38.14\,\mathrm{dB}),\\
A_{CM}&\approx-\frac{1}{2R_{SS}g_{m9}}=-1.4482\times10^{-4},\\
\mathrm{CMRR}&\approx20\log_{10}\left|\frac{80.731}{-1.4482\times10^{-4}}\right|
=114.92\,\mathrm{dB}.
\end{align*}
\begin{center}
\small
\begin{tabular}{@{}lcc@{}}
\toprule
Specification & Requirement & Analytical estimate\\
\midrule
Slew rate, $I_{\mathrm{tail}}/C_L$ & $>7\Vus$ & $12\Vus$\\
GBW, $g_{m1}/(2\pi C_L)$ & $>8\MHz$ & $12.73\MHz$\\
CMRR & $>65$ dB & $114.92$ dB\\
Common-mode range width & Nonzero & $0.7432\V$\\
Dominant pole, $1/(2\pi R_{\mathrm{out}}C_L)$ & Record estimate & $157.71$ kHz\\
\bottomrule
\end{tabular}
\end{center}

\end{enumerate}

\newpage
\textbf{Model validity and comparison with simulation results.}

The calculated dimensions satisfy the requirements within the stated
first-order model. The Lab 1 parameters were extracted at one geometry
and bias condition. Body effect, mobility degradation, finite
drain-voltage-dependent currents, internal poles, and parasitic
capacitances are not fully represented by these estimates. Pole and
GBW calculations assume the $10$ pF load dominates the output
capacitance and higher-frequency internal poles are negligible.
The Lab 1 report provides $\theta$ but no numerical body-effect
parameters. Threshold shifts caused by the elevated NMOS source
voltages and PMOS cascode source voltages therefore remain
unquantified. The computed CMRR assumes matched devices and ideal
AC-grounded cascode biases, excluding mismatch effects.

The Lab 5 report illustrates the difference between analytical
sizing and simulation-refined dimensions. Its simulated $120\uA$
simple mirror uses $W=19.4\um$,
$L=0.270\um$ for approximately $0.150\V$ compliance, whereas
the first-order square-law calculation gives $W=13.36\um$.
The final cascode mirror uses $L=0.320\um$ to increase output
resistance. These dimensions characterize the refined Lab 5 circuits;
they do not establish a new device-parameter set for the Lab 6 designs.
The present design retains $L=0.270\um$ to use the corresponding
Lab 1 parameter set without assigning an unmeasured $\lambda$
to $0.320\um$ devices.

A common-mode DC sweep determines the interval in which all active
devices remain in saturation. The widths, reference current, and
$V_{BN}/V_{BP}$ control the bias point and provide the variables
for design refinement. Separate differential-mode
and common-mode AC excitations determine the corresponding gains,
CMRR, GBW, and dominant pole. Slew rate requires a large-signal
transient analysis because it lies outside the linearized AC model.
Saturation and performance margins account for the effects omitted
from the first-order design.

\vspace{12pt}
\textbf{References.}
\begin{enumerate}[label={[\arabic*]},leftmargin=0.3in,
topsep=4pt,itemsep=5pt,parsep=0pt]
\item S. Palermo, \emph{ECEN 474/704 Lab 6: Differential Pairs},
course laboratory handout, Figs.~6-6 and 6-8 and Prelab Questions 1--3.
\item Anonymous contributor, \emph{ECEN 704 Lab 1 Report}, Table I:
single-finger $5/0.27\um$ NFET and PFET parameter rows.
\item Anonymous contributor, \emph{ECEN 704 Lab 5 Report},
Sec.~III and Table 3: final mirror sizes and simulation-based refinement.
\end{enumerate}

\end{document}
